% ---- begin paper/tables/a12_pscs.tex
\begin{table}[htbp]\centering
\caption{Caregiver Measures II: Parenting Sense of Competence}
\label{tab:a_pscs}
\footnotesize\renewcommand{\arraystretch}{0.92}
\begin{threeparttable}
\begin{tabular}{@{}p{68pt}p{212pt}p{150pt}@{}}
\toprule
 & Question, verbatim & Response scale \\
 & (1) & (2) \\
\midrule
\texttt{pscs\_p1} & Los problemas de cuidar a un(a) adolescente o joven son fáciles de resolver una vez que sabes cómo tus acciones lo pueden afectar & 1 = Muy en desacuerdo; 2 = Moderadamente en desacuerdo; 3 = Ni de acuerdo ni en desacuerdo; 4 = Moderadamente de acuerdo; 5 = Muy de acuerdo \\
\addlinespace[1pt]
\texttt{pscs\_p2} & Cumplo mis expectativas respecto al dominio que tengo en el cuidado del adolescente o joven & \textit{Same scale} \\
\addlinespace[1pt]
\texttt{pscs\_p3} & Yo sería un buen ejemplo a seguir para que una madre o padre aprendan lo necesario para ser buenos padres & \textit{Same scale} \\
\addlinespace[1pt]
\texttt{pscs\_p4} & Ser cuidador es manejable, y los problemas se resuelven fácilmente & \textit{Same scale} \\
\addlinespace[1pt]
\texttt{pscs\_p5} & Si alguien puede encontrar la respuesta a lo que le molesta al adolescente o joven, ese/esa soy yo & \textit{Same scale} \\
\addlinespace[1pt]
\texttt{pscs\_p6} & Un problema difícil de ser cuidador(a) es no saber si estás haciendo un buen trabajo & \textit{Same scale} \\
\addlinespace[1pt]
\texttt{pscs\_p7} & Teniendo en cuenta el tiempo que he sido un(a) cuidador(a), me siento muy cómodo(a) en este rol & \textit{Same scale} \\
\addlinespace[1pt]
\texttt{pscs\_p8} & Creo que tengo todas las habilidades necesarias para ser un(a) buen(a) cuidador(a) para el adolescente o joven & \textit{Same scale} \\
\addlinespace[1pt]
\texttt{pscs\_p9} & A pesar de que ser cuidador(a) puede ser gratificante, me siento frustrado(a) por como es el adolescente o joven en este momento & \textit{Same scale} \\
\addlinespace[1pt]
\texttt{pscs\_p10} & No sé por qué, pero a veces cuando se supone que debo tener el control, más bien siento que estoy siendo manipulado(a) & \textit{Same scale} \\
\addlinespace[1pt]
\texttt{pscs\_p11} & Mi padre/madre estaba mejor preparado(a) que yo para ser un(a) buen(a) cuidador(a) & \textit{Same scale} \\
\addlinespace[1pt]
\texttt{pscs\_p12} & A veces siento como si no estuviera logrando terminar nada & \textit{Same scale} \\
\addlinespace[1pt]
\texttt{pscs\_p13} & Voy a la cama de la misma manera que me despierto por la mañana, sintiendo que no he logrado mucho & \textit{Same scale} \\
\addlinespace[1pt]
\texttt{pscs\_p14} & Mis talentos e intereses están en otras áreas, no en ser un/a cuidador/a & \textit{Same scale} \\
\addlinespace[1pt]
\texttt{pscs\_p15} & Si ser cuidador(a) de un adolescente o joven fuera más interesante, estaría más motivado(a) a hacer un mejor trabajo como cuidador(a) & \textit{Same scale} \\
\addlinespace[1pt]
\texttt{pscs\_p16} & Ser cuidador/a me pone tenso/a y ansioso/a & \textit{Same scale} \\
\addlinespace[1pt]
\texttt{pscs\_p17} & Ser un/a buen/a cuidador/a es una recompensa en sí mismo & \textit{Same scale} \\
\bottomrule
\end{tabular}
\begin{wpnotes}
\textit{Notes:} Question wording and response labels are reproduced verbatim in the original Spanish of the questionnaire, from the variable and value labels of the public microdata; wording the microdata truncates at 80 characters is marked with an ellipsis. Parenting Sense of Competence Scale, answered by the main caregiver; the analysis uses the survey's total score, standardized, with higher values meaning more competence.
\end{wpnotes}
\end{threeparttable}
\end{table}
% ---- end paper/tables/a12_pscs.tex
